% main_prob_lecture5_metropolis.tex
\documentclass[aspectratio=169]{beamer}

% --- Language & encoding ---
\usepackage[utf8]{inputenc}
\usepackage[T1]{fontenc}
\usepackage[english,french]{babel}
\usepackage{lmodern}
\usepackage{microtype}

% --- Math & tables ---
\usepackage{amsmath,amssymb}
\usepackage{booktabs}
\usepackage{graphicx}
\usepackage{changepage}

% --- Metropolis Theme ---
\usetheme[numbering=fraction,progressbar=frametitle]{metropolis}
\setbeamertemplate{frame numbering}[fraction]

% --- Colors (same palette as prior lectures) ---
\definecolor{BootcampBlueDark}{HTML}{1A3C68}
\definecolor{TextBlack}{HTML}{000000}
\definecolor{AccentGold}{HTML}{DAA520}
\definecolor{Crimson}{HTML}{990000}
\definecolor{MatrixGreen}{HTML}{228B22}
\definecolor{MatrixRed}{HTML}{B22222}

\setbeamercolor{block title example}{bg=BootcampBlueDark!90!black,fg=white}
\setbeamercolor{block body example}{bg=BootcampBlueDark!5!white,fg=black}

% === Thick frame title band ===
\setbeamercolor{frametitle}{bg=BootcampBlueDark,fg=white}
\setbeamerfont{frametitle}{series=\bfseries,size=\large}
\setbeamertemplate{frametitle}{
  \nointerlineskip
  \begin{beamercolorbox}[wd=\paperwidth,ht=3.2ex,dp=1.4ex,leftskip=1em,rightskip=1em]{frametitle}
    \usebeamerfont{frametitle}\insertframetitle
  \end{beamercolorbox}
}

% === Custom Definition Block ===
\newenvironment{definitionblock}[1]{%
  \par\vspace{0.4em}%
  \begin{beamercolorbox}[
      wd=\textwidth,
      sep=0.4em,
      leftskip=0.3em,
      rightskip=0.6em,
      rounded=false,
      shadow=false
    ]{block title example}%
    \raisebox{0.15em}{\usebeamerfont{block title example}#1}%
  \end{beamercolorbox}%
  \vspace{-0.4em}%
  \begin{beamercolorbox}[
      wd=\textwidth,
      sep=1em,
      leftskip=0.6em,
      rightskip=0.6em,
      rounded=false,
      shadow=false
    ]{block body example}%
    \usebeamerfont{block body example}%
}{%
  \end{beamercolorbox}%
  \par\vspace{0.6em}%
}

% === Global text formatting ===
\setbeamercolor{normal text}{fg=TextBlack,bg=white}
\setbeamercolor{title}{fg=TextBlack}
\setbeamercolor{structure}{fg=BootcampBlueDark}
\setbeamercolor{progress bar}{fg=BootcampBlueDark,bg=gray!30}

% --- Course Info ---
\title{Probability ---Conditional Laws, Expectation \& Martingales }
\author{Sauldubois Nathan}
\institute{NYU Tandon School, Finance and Risk Engineering}
\date{\today}
% \titlegraphic{\includegraphics[height=1.2cm]{your_logo.pdf}}

\metroset{
  sectionpage=progressbar,
  subsectionpage=progressbar
}

% --- Section & subsection transitions (plain) ---
\AtBeginSection{
  \frame[plain]{\sectionpage}
}
\AtBeginSubsection{
  \frame[plain]{\subsectionpage}
}

\begin{document}

\maketitle

% ===== Outline =====
\begin{frame}{Outline}
  \tableofcontents[hideallsubsections]
\end{frame}


% ======================================
% SECTION 2 — Conditional Expectation : loi conditionel, proba conditionel etc programmation dynamique, chaine de markov 
% ======================================

% ======================================
% SECTION — Conditional Expectation & Information
% ======================================
\section{Conditional Expectation \& Information}

\begin{frame}{Outline}
  \tableofcontents[currentsection, hideothersubsections]
\end{frame}

% ======================================
\subsection{Sigma-Algebras \& Information}

\begin{frame}{Information Sets}

\begin{definitionblock}{\(\sigma\)-Algebra}
A \emph{\(\sigma\)-algebra} \(\mathcal{G}\subseteq\mathcal{F}\) is a collection of events such that:
\begin{itemize}
  \item \(\Omega\in\mathcal{G}\), \pause
  \item if \(A\in\mathcal{G}\), then \(A^c\in\mathcal{G}\),\pause
  \item if \((A_n)\subset\mathcal{G}\), then \(\bigcup_n A_n\in\mathcal{G}\).\pause
\end{itemize}
\end{definitionblock}

\medskip

\begin{itemize}
  \item A \(\sigma\)-algebra represents the \emph{information available}.\pause
  \item More events \(\Rightarrow\) more information.
\end{itemize}

\end{frame}

\begin{frame}{Example: Partial Information on a Finite Space}

Consider a dice roll:
\[
\Omega=\{1,2,3,4,5,6\}.
\]
\pause
\medskip

\textbf{Full information:}
\[
\mathcal{F}=\mathcal{P}(\Omega)
\quad\text{(all subsets)}.
\]\pause

\medskip

\textbf{Partial information:}
suppose we only observe whether the result is even or odd.
\[
\mathcal{G}
=
\sigma(\{2,4,6\})
=
\{\emptyset,\{2,4,6\},\{1,3,5\},\Omega\}.
\]
\pause
\medskip

\begin{itemize}
  \item We cannot distinguish 2 from 4 or 6.
  \item Events not in \(\mathcal{G}\) are \emph{unobservable}.
\end{itemize}

\end{frame}

\begin{frame}{Example: Information Generated by a Random Variable}

Let \(X:\Omega\to\mathbb{R}\) be a random variable.
\pause
\medskip

The \(\sigma\)-algebra generated by \(X\) is
\[
\sigma(X)
=
\{X^{-1}(B): B\in\mathcal{B}(\mathbb{R})\}.
\]\pause

\medskip

\begin{itemize}
  \item \(\sigma(X)\) contains all events that can be described using \(X\).
  \item Knowing \(\sigma(X)\) is equivalent to knowing the value of \(X\).
\end{itemize}
\pause
\medskip

\textbf{Example:}
\begin{itemize}
  \item \(X=\mathbf{1}_A\)  \(\Rightarrow\)
  \(\sigma(X)=\{\emptyset,A,A^c,\Omega\}\).
\end{itemize}
\pause
\medskip

\textbf{Interpretation:}
conditional expectation given \(\sigma(X)\) means
\emph{using only the information carried by \(X\)}.

\end{frame}


% --------------------------------------

\begin{frame}{Filtrations: Information Over Time}

\begin{definitionblock}{Filtration}
A \emph{filtration} is an increasing family of \(\sigma\)-algebras
\[
(\mathcal{F}_t)_{t\ge0}
\quad\text{with}\quad
\mathcal{F}_s\subseteq\mathcal{F}_t
\ \text{for } s\le t.
\]
\end{definitionblock}\pause

\begin{itemize}
  \item \(\mathcal{F}_t\): information available up to time \(t\).
  \item Natural framework for time series and stochastic processes.
\end{itemize}\pause

\medskip

\textbf{Example:}
\begin{itemize}
  \item \(\mathcal{F}_t = \sigma(X_1,\dots,X_t)\) generated by observations up to time \(t\).
\end{itemize}

\end{frame}




% ======================================
\subsection{Conditional Distributions}

\begin{frame}{Conditional Distributions}

\textbf{Discrete case.}

Let \((X,Y)\) be discrete random variables.\pause
For \(y\) such that \(\mathbb{P}(Y=y)>0\),\pause
\[
\mathbb{P}(X=x\mid Y=y)
=
\frac{\mathbb{P}(X=x,Y=y)}{\mathbb{P}(Y=y)}.
\]\pause

\medskip

\textbf{Continuous case.}

If \((X,Y)\) has joint density \(f_{X,Y}\),\pause
the conditional density of \(X\mid Y=y\) is\pause
\[
f_{X\mid Y=y}(x)
=
\frac{f_{X,Y}(x,y)}{f_Y(y)},
\quad \text{if } f_Y(y)>0.
\]\pause

\medskip

If \(X\) and \(Y\) are independent, conditioning has no effect.

\end{frame}

\begin{frame}{Example: Discrete Conditional Distribution}

Let \((X,Y)\) take values in \(\{0,1\}\times\{0,1\}\) with joint probabilities:\pause
\[
\begin{array}{c|cc}
 & Y=0 & Y=1 \\ \hline
X=0 & 0.3 & 0.1 \\
X=1 & 0.2 & 0.4
\end{array}
\]
\pause
\medskip

Compute the marginal of \(Y\):
\[
\mathbb{P}(Y=1)=0.1+0.4=0.5.
\]
\pause
\medskip

Then the conditional distribution of \(X\mid Y=1\) is:
\[
\mathbb{P}(X=1\mid Y=1)=\frac{0.4}{0.5}=0.8,
\quad
\mathbb{P}(X=0\mid Y=1)=0.2.
\]\pause

\medskip

\textbf{Interpretation:}
knowing \(Y=1\) strongly increases the likelihood of \(X=1\).

\end{frame}

\begin{frame}{Why Conditional Distributions Matter (Finance Example)}
\pause
\textbf{Market regime and asset returns.}
\pause
\begin{itemize}
  \item \(R\): daily return of a stock,\pause
  \item \(M\): market regime indicator,\pause
  $
  M =
  \begin{cases}
  1 & \text{bull market},\\
  0 & \text{bear market}.
  \end{cases}
  $
\end{itemize}
\pause
\medskip

Empirically, the distribution of returns depends on the regime:\pause
\[
R \mid (M=1) \sim \mathcal{N}(\mu_{\text{bull}},\sigma_{\text{bull}}^2),
\qquad
R \mid (M=0) \sim \mathcal{N}(\mu_{\text{bear}},\sigma_{\text{bear}}^2),
\]
with typically \(\mu_{\text{bull}}>\mu_{\text{bear}}\) and
\(\sigma_{\text{bear}}>\sigma_{\text{bull}}\).
\pause
\medskip

\textbf{Key message:}
\begin{itemize}
  \item The unconditional distribution of \(R\) mixes very different risks.\pause
  \item Risk management and forecasting rely on
  \[
  \textbf{conditional distributions } \; \mathcal{L}(R \mid M).
  \]\pause
\end{itemize}


\end{frame}

\begin{frame}{Conditional Probability: The General Idea}

\begin{definitionblock}{Conditional Probability}
Let \(A,B\) be events with \(\mathbb{P}(B)>0\).\pause
The \emph{conditional probability} of \(A\) given \(B\) is
\[
\mathbb{P}(A\mid B)
=
\frac{\mathbb{P}(A\cap B)}{\mathbb{P}(B)}.
\vspace{-5mm}
\]\pause
\vspace{-5mm}
\end{definitionblock}
\vspace{-3mm}

\textbf{Interpretation:}
\begin{itemize}
  \item We restrict attention to the scenario where \(B\) has occurred.\pause
  \item Probabilities are \emph{renormalized} on \(B\).\pause
\end{itemize}
\vspace{-3mm}

\textbf{From events to random variables:}
\begin{itemize}
  \item Conditioning on an event \(B\)  
  \(\longrightarrow\) \(\mathbb{P}(\cdot\mid B)\)\pause
  \item Conditioning on a random variable \(Y\)  
  \(\longrightarrow\) conditional laws \(\mathcal{L}(X\mid Y=y)\)\pause
\end{itemize}

\medskip

\textbf{Key principle:}
\begin{quote}
Conditioning means updating probabilities after observing information.
\end{quote}

\end{frame}


\begin{frame}{Example: Continuous Conditional Density}

Let \((X,Y)\) have joint density
\[
f_{X,Y}(x,y)=
\begin{cases}
2 & \text{if } 0<x<y<1,\\
0 & \text{otherwise.}
\end{cases}
\vspace{-5mm}
\]\pause

\medskip

First compute the marginal density of \(Y\):
\[
f_Y(y)=\int_0^y 2\,dx = 2y,
\quad 0<y<1.
\vspace{-5mm}
\]
\pause
\medskip

The conditional density of \(X\mid Y=y\) is
\[
f_{X\mid Y=y}(x)=\frac{2}{2y}=\frac{1}{y},
\quad 0<x<y.
\vspace{-5mm}
\]
\pause
\medskip

\textbf{Conclusion:}
\[
X\mid Y=y \sim \mathrm{Uniform}(0,y).
\]

\end{frame}

\begin{frame}{Example: Independence and Conditioning}

Let \(X\sim \mathcal{N}(0,1)\) and \(Y\sim \mathrm{Exp}(1)\),
independent.
\pause
\medskip

Then the joint density factorizes:
\[
f_{X,Y}(x,y)=f_X(x)f_Y(y).
\]
\pause
\medskip

Hence, for all \(y>0\),
\[
f_{X\mid Y=y}(x)=f_X(x).
\]
\pause
\medskip

\textbf{Interpretation:}
\begin{itemize}
  \item Knowing \(Y\) gives no information about \(X\).\pause
  \item Conditioning does not change the distribution.
\end{itemize}

\end{frame}

% ======================================
\subsection{Conditional Expectation}

\begin{frame}{Conditional Expectation: Motivation}

\begin{itemize}
  \item What is the \emph{best prediction} of a random variable \(X\),
  given partial information?\pause
  \item Conditional expectation answers this question.\pause
\end{itemize}

\medskip

\textbf{Goal:}
predict \(X\) using only information contained in:
\begin{itemize}
  \item a random variable \(Y\), or\pause
  \item a \(\sigma\)-algebra \(\mathcal{G}\).
\end{itemize}

\end{frame}

% --------------------------------------

\begin{frame}{\(\mathbb{E}[X\mid Y]\): Concrete Definition}

\textbf{Discrete case.}

If \((X,Y)\) is discrete,
\[
\mathbb{E}[X\mid Y]
=
\sum_x x\,\mathbb{P}(X=x\mid Y).
\]\pause

\medskip

\textbf{Continuous case.}

If \((X,Y)\) has joint density,
\[
\mathbb{E}[X\mid Y]
=
\int x\, f_{X\mid Y}(x\mid Y)\,dx.
\]\pause

\medskip

\begin{itemize}
  \item \(\mathbb{E}[X\mid Y]\) is a \emph{function of \(Y\)}.\pause
  \item Hence it is itself a random variable.
\end{itemize}

\end{frame}


\begin{frame}{Conditional Expectation Given $\sigma(Y)$ (Test Functions Formulation)}

\begin{definitionblock}{Conditional Expectation w.r.t. $\sigma(Y)$}
Let $X\in L^1(\Omega,\mathcal F,\mathbb P)$ and let $\mathcal G=\sigma(Y)$ for some random variable $Y$.
The conditional expectation $\mathbb{E}[X\mid\sigma(Y)]$ is the unique $\sigma(Y)$-measurable random variable $m(Y)$ such that
\[
\mathbb{E}\!\big[m(Y)\,\varphi(Y)\big] \;=\; \mathbb{E}\!\big[X\,\varphi(Y)\big]
\quad
\text{for all bounded measurable functions } \varphi.
\]
\end{definitionblock}

\pause
\medskip

\begin{itemize}
  \item Here the \emph{test functions} are exactly functions of $Y$: $Z=\varphi(Y)$.
  \pause
  \item This is equivalent to testing only indicators $\varphi=\mathbf 1_{B}$ (hence events $\{Y\in B\}$).
  \pause
  \item Emphasizes $\mathbb{E}[X\mid\sigma(Y)]$ as the best $\sigma(Y)$-based prediction (a projection onto functions of $Y$).
\end{itemize}

\end{frame}


% --------------------------------------

\begin{frame}{Example 1: Discrete Conditional Expectation}

Let \((X,Y)\) be discrete with joint distribution:
\[
\begin{array}{c|cc}
 & Y=0 & Y=1 \\ \hline
X=0 & 0.3 & 0.1 \\
X=1 & 0.2 & 0.4
\end{array}
\vspace{-5mm}
\]\pause

\medskip

Compute \(\mathbb{E}[X\mid Y]\).
\pause
\medskip

For \(Y=1\):
\[
\mathbb{E}[X\mid Y=1]
=
0\cdot 0.2 + 1\cdot 0.8 = 0.8.
\vspace{-5mm}
\]
\pause
For \(Y=0\):
\[
\mathbb{E}[X\mid Y=0]
=
0\cdot \tfrac{3}{5} + 1\cdot \tfrac{2}{5} = \tfrac{2}{5}.
\vspace{-5mm}
\]
\pause

Hence
\[
\mathbb{E}[X\mid Y]
=
0.8\,\mathbf{1}_{\{Y=1\}}
+
\tfrac{2}{5}\,\mathbf{1}_{\{Y=0\}}.
\]

\end{frame}

\begin{frame}{Example 2: Continuous Conditional Expectation}

Let \((X,Y)\) have joint density
\[
f_{X,Y}(x,y)=
\begin{cases}
2 & \text{if } 0<x<y<1,\\
0 & \text{otherwise.}
\end{cases}
\]\pause

\medskip

We have seen that
\[
X\mid Y=y \sim \mathrm{Uniform}(0,y).
\vspace{-5mm}
\]\pause

\medskip

Therefore
\[
\mathbb{E}[X\mid Y=y]
=
\int_0^y x\,\frac{1}{y}\,dx
=
\frac{y}{2}.
\vspace{-5mm}
\]\pause

\medskip

Conclusion:
\[
\mathbb{E}[X\mid Y]=\frac{Y}{2}.
\]

\end{frame}

\begin{frame}{Example 3: Independence}

Let \(X\sim\mathcal{N}(0,1)\),
\(Y\sim\mathrm{Exp}(1)\),
independent.
\pause
\medskip

Since \(X\) is independent of \(Y\),
\[
\mathbb{E}[X\mid Y]=\mathbb{E}[X]=0.
\vspace{-5mm}
\]
\pause

More generally, for any integrable function \(\varphi\),
\[
\mathbb{E}[\varphi(X)\mid Y]
=
\mathbb{E}[\varphi(X)].
\vspace{-5mm}
\]\pause

\medskip

\textbf{Key point:}
conditioning on independent information does nothing.

\end{frame}

\begin{frame}{Example 4: Projection onto a Random Variable}

Let \(X,Y\in L^2\).

We look for
\[
\mathbb{E}[X\mid Y]
=
a+bY
\quad\text{(best linear predictor)}.
\]\pause
The orthogonality conditions give:
\[
b=\frac{\mathrm{Cov}(X,Y)}{\mathrm{Var}(Y)},
\qquad
a=\mathbb{E}[X]-b\,\mathbb{E}[Y].
\]\pause

Hence
\[
\mathbb{E}[X\mid Y]
=
\mathbb{E}[X]
+
\frac{\mathrm{Cov}(X,Y)}{\mathrm{Var}(Y)}
\big(Y-\mathbb{E}[Y]\big).
\]
\pause
\medskip

This is exact when \((X,Y)\) is Gaussian.

\end{frame}

\begin{frame}{Example 5: Conditional Expectation for Gaussian Vectors}

Let
\[
\begin{pmatrix} X \\ Y \end{pmatrix}
\sim
\mathcal{N}\!\left(
\begin{pmatrix} \mu_X \\ \mu_Y \end{pmatrix},
\begin{pmatrix}
\sigma_X^2 & \rho\sigma_X\sigma_Y \\
\rho\sigma_X\sigma_Y & \sigma_Y^2
\end{pmatrix}
\right).
\]\pause

\medskip

Then
\[
\mathbb{E}[Y\mid X]
=
\mu_Y
+
\rho\frac{\sigma_Y}{\sigma_X}(X-\mu_X).
\]
\pause
\medskip

\textbf{Key properties:}
\begin{itemize}
  \item affine function of \(X\),
  \item optimal predictor in mean-square sense.
\end{itemize}

\end{frame}

\begin{frame}{Example 6: Financial Interpretation}

Let \(R\) be tomorrow's return and \(S\) a signal observed today.
\pause
Assume \((R,S)\) is jointly Gaussian.
\pause
\medskip

Then the best forecast of \(R\) given \(S\) is:
\[
\mathbb{E}[R\mid S]
=
\mathbb{E}[R]
+
\frac{\mathrm{Cov}(R,S)}{\mathrm{Var}(S)}
\big(S-\mathbb{E}[S]\big).
\]\pause

\medskip

\textbf{Interpretation:}
\begin{itemize}
  \item linear factor model,\pause
  \item foundation of predictive regressions,\pause
  \item links information to expected return.
\end{itemize}

\end{frame}

% ======================================
\subsection{Properties}

\begin{frame}{Algebraic Properties}

Let \(X,X'\in L^1\), \(Y\) a random variable, \(\mathcal{G}=\sigma(Y)\).
\pause
\begin{itemize}
  \item \textbf{Linearity:}
  \[
  \mathbb{E}[aX+X'\mid\mathcal{G}]
  =
  a\,\mathbb{E}[X\mid\mathcal{G}]
  +\mathbb{E}[X'\mid\mathcal{G}].
  \]
\pause
  \item \textbf{Constants:}
  if \(c\in\mathbb{R}\),
  \[
  \mathbb{E}[c\mid\mathcal{G}]=c.
  \]
\pause
  \item \textbf{Taking out what is known:}
  if \(Z\) is \(\mathcal{G}\)-measurable,
  \[
  \mathbb{E}[ZX\mid\mathcal{G}]
  =
  Z\,\mathbb{E}[X\mid\mathcal{G}].
  \]\pause
\end{itemize}

\medskip

These properties make conditional expectation a linear operator.

\end{frame}

\begin{frame}{Order and Monotonicity}

Let \(X,X'\in L^1\), \(\mathcal{G}\subseteq\mathcal{F}\).

\begin{itemize}
  \item \textbf{Positivity:}
  if \(X\ge0\) a.s., then
  \[
  \mathbb{E}[X\mid\mathcal{G}]\ge0 \quad \text{a.s.}
  \]\pause

  \item \textbf{Monotonicity:}
  if \(X\le X'\) a.s., then
  \[
  \mathbb{E}[X\mid\mathcal{G}]
  \le
  \mathbb{E}[X'\mid\mathcal{G}]
  \quad \text{a.s.}
  \]\pause

  \item \textbf{Bounds preserved:}
  if \(a\le X\le b\),
  \[
  a \le \mathbb{E}[X\mid\mathcal{G}] \le b.
  \]\pause
\end{itemize}

\medskip

Conditional expectation preserves order and inequalities.

\end{frame}

\begin{frame}{Conditional Jensen Inequality}

Let \(\varphi:\mathbb{R}\to\mathbb{R}\) be convex and \(X\in L^1\).

\begin{definitionblock}{Conditional Jensen Inequality}
\[
\varphi\big(\mathbb{E}[X\mid\mathcal{G}]\big)
\le
\mathbb{E}[\varphi(X)\mid\mathcal{G}]
\quad \text{a.s.}
\]\pause
\end{definitionblock}

\medskip

\begin{itemize}
  \item Generalizes Jensen’s inequality.\pause
  \item Fundamental in risk measures and utility theory.\pause
  \item Equality holds if \(X\) is \(\mathcal{G}\)-measurable.
\end{itemize}

\end{frame}

\begin{frame}{Tower Property and Filtrations}

Let \(\mathcal{H}\subseteq\mathcal{G}\subseteq\mathcal{F}\).\pause

\begin{definitionblock}{Tower Property}
\[
\mathbb{E}\!\left[\mathbb{E}[X\mid\mathcal{G}]\mid\mathcal{H}\right]
=
\mathbb{E}[X\mid\mathcal{H}].
\]\pause
\end{definitionblock}

\medskip

\begin{itemize}
  \item Conditioning step by step does not add information.\pause
  \item Essential for dynamic models and martingales.
\end{itemize}

\end{frame}

\begin{frame}{Contraction and Convergence Properties}

Let \(p\ge1\) and \(X\in L^p\).

\begin{itemize}
  \item \textbf{Contraction in \(L^p\):}
  \[
  \|\mathbb{E}[X\mid\mathcal{G}]\|_p
  \le
  \|X\|_p.
  \]\pause

  \item \textbf{Continuity in \(L^p\):}
  if \(X_n\xrightarrow{L^p} X\), then
  \[
  \mathbb{E}[X_n\mid\mathcal{G}]
  \xrightarrow{L^p}
  \mathbb{E}[X\mid\mathcal{G}].
  \]\pause

  \item \textbf{Almost sure convergence:}
  if \(X_n\to X\) a.s.\ and \(|X_n|\le Y\in L^1\),
  \[
  \mathbb{E}[X_n\mid\mathcal{G}]
  \to
  \mathbb{E}[X\mid\mathcal{G}]
  \quad \text{a.s.}
  \]
\end{itemize}

\end{frame}

\begin{frame}{Variance Decomposition}

If \(X\in L^2\), then:
\[
\mathrm{Var}(X)
=
\mathbb{E}\!\left[\mathrm{Var}(X\mid\mathcal{G})\right]
+
\mathrm{Var}\!\left(\mathbb{E}[X\mid\mathcal{G}]\right).
\]\pause

\medskip

\begin{itemize}
  \item Total risk = unexplained risk + explained risk.\pause
  \item Core result in statistics and finance.
\end{itemize}

\end{frame}


% ======================================
\subsection{Projection Interpretation}

\begin{frame}{Conditional Expectation as Projection}

Let \(X,Y\in L^2\).
Define
\[
L^2(Y)=\{g(Y):\ \mathbb{E}[g(Y)^2]<\infty\}.
\]\pause

\medskip

\begin{definitionblock}{Projection Theorem}
\[
\mathbb{E}[X\mid Y]
=
\arg\min_{Z\in L^2(Y)}
\mathbb{E}\big[(X-Z)^2\big].
\]\pause
\end{definitionblock}

\begin{itemize}
  \item Best mean-square predictor of \(X\) using \(Y\).\pause
  \item Error is orthogonal to all functions of \(Y\).
\end{itemize}

\end{frame}

\begin{frame}{Geometric Interpretation in \texorpdfstring{$L^2$}{L2}}

\centering
\begin{tikzpicture}[scale=1.05]

% Axes / space
\draw[thick] (-0.5,0) -- (5,0);
\node at (5.2,0) {\(L^2(Y)\)};

% Subspace
\draw[very thick, BootcampBlueDark] (0,0) -- (4.5,0);
\node[above] at (3.2,0.15) {\(\text{functions of } Y\)};

% X vector
\draw[->, thick] (0,0) -- (3,2.5);
\node[above] at (3,2.5) {\(X\)};

% Projection
\draw[dashed] (3,2.5) -- (3,0);
\draw[->, thick, AccentGold] (0,0) -- (3,0);
\node[below] at (3,0) {\(\mathbb{E}[X\mid Y]\)};

% Error
\draw[->, thick, Crimson] (3,0) -- (3,2.5);
\node[right] at (3.05,1.25) {\(X-\mathbb{E}[X\mid Y]\)};

% Orthogonality symbol
\node at (3.1,0.25) {\(\perp\)};

\end{tikzpicture}

\medskip

\begin{itemize}
  \item \(\mathbb{E}[X\mid Y]\) is the orthogonal projection of \(X\) onto \(L^2(Y)\).\pause
  \item The error is orthogonal to all functions of \(Y\).\pause
  \item Minimizes the mean squared error.
\end{itemize}

\end{frame}

% ======================================


% ======================================
\subsection{Conditioning and Gaussian Vectors}

\begin{frame}{Gaussian Vectors: A Key Property}

\begin{definitionblock}{Stability under Conditioning}
If a random vector \((X,Y)\) is jointly Gaussian, then:
\begin{itemize}
  \item the conditional distribution of \(Y\mid X\) is Gaussian;
  \item the conditional expectation \(\mathbb{E}[Y\mid X]\) is an \emph{affine} function of \(X\).
\end{itemize}
\end{definitionblock}\pause

\medskip

This property is specific to the Gaussian family and does not hold in general.

\end{frame}

% --------------------------------------

\begin{frame}{Conditional Distribution of a Gaussian Vector}

Let
\[
\begin{pmatrix} X \\ Y \end{pmatrix}
\sim
\mathcal{N}\!\left(
\begin{pmatrix} \mu_X \\ \mu_Y \end{pmatrix},
\begin{pmatrix}
\Sigma_{XX} & \Sigma_{XY} \\
\Sigma_{YX} & \Sigma_{YY}
\end{pmatrix}
\right),
\]
with \(\Sigma_{XX}\) invertible.\pause

\medskip

Then the conditional distribution of \(Y\mid X\) is Gaussian:
\[
Y\mid X
\sim
\mathcal{N}\!\big(
\mu_{Y\mid X},
\Sigma_{Y\mid X}
\big),
\]
where
\[
\mu_{Y\mid X}
=
\mu_Y
+
\Sigma_{YX}\Sigma_{XX}^{-1}(X-\mu_X),
\]
\[
\Sigma_{Y\mid X}
=
\Sigma_{YY}
-
\Sigma_{YX}\Sigma_{XX}^{-1}\Sigma_{XY}.
\]

\end{frame}

% --------------------------------------

\begin{frame}{Interpretation and Applications}

\begin{itemize}
  \item \textbf{Linear prediction:}
  \[
  \mathbb{E}[Y\mid X]
  =
  \mu_Y
  +
  \Sigma_{YX}\Sigma_{XX}^{-1}(X-\mu_X).
  \]
  \item Conditional variance does not depend on the realized value of \(X\).
  \item Conditioning reduces uncertainty:
  \[
  \Sigma_{Y\mid X} \le \Sigma_{YY}.
  \]
\end{itemize}

\medskip

\textbf{Applications:}
\begin{itemize}
  \item Linear regression.
  \item Kalman filtering.
  \item Portfolio forecasting and risk management.
\end{itemize}

\end{frame}

% ======================================
% SECTION 3 — (Sub/Super)Martingales
% ======================================
% ======================================
\section{Martingales, Submartingales, Supermartingales}

\begin{frame}{Outline}
  \tableofcontents[currentsection, hideothersubsections]
\end{frame}

% ======================================
\subsection{Definitions}

\begin{frame}{Martingales and Natural Filtration}

Let \((X_t)_{t\ge0}\) be a stochastic process.\pause

\medskip

\begin{itemize}
  \item The \textbf{natural filtration} of \(X\) is
  \[
  \mathcal{F}_t = \sigma(X_0,\dots,X_t),
  \]\pause
  representing all information revealed by the process up to time \(t\).
  \item We assume \(X_t\in L^1\) for all \(t\).
\end{itemize}

\end{frame}

% --------------------------------------

\begin{frame}{Fair Game Interpretation}

\begin{definitionblock}{Martingale / Submartingale / Supermartingale}
The process \((X_t)\) is:
\[
\mathbb{E}[X_{t+1}\mid \mathcal{F}_t] =
\begin{cases}
X_t & \text{martingale} \\[1mm]
\ge X_t & \text{submartingale} \\[1mm]
\le X_t & \text{supermartingale}
\end{cases}
\]
\end{definitionblock}

\medskip

\begin{itemize}
  \item Martingale: \emph{fair game}, no predictable drift.\pause
  \item Submartingale: favorable trend.\pause
  \item Supermartingale: unfavorable trend.
\end{itemize}

\end{frame}

% ======================================
\subsection{Canonical Examples}

\begin{frame}{Example 1: Sums of i.i.d.\ Variables}

Let \(X_1,X_2,\dots\) be i.i.d.\ with
\(\mathbb{E}[X_1]=0\).
\pause
Define
\[
S_t = \sum_{k=1}^t X_k,
\qquad
\mathcal{F}_t=\sigma(X_1,\dots,X_t).
\]
\pause
\medskip

Then
\[
\mathbb{E}[S_{t+1}\mid\mathcal{F}_t]
=
S_t + \mathbb{E}[X_{t+1} \vert \mathcal{F}_t ]
=
S_t.
\]

\medskip

\textbf{Conclusion:}
\((S_t)\) is a martingale.

\end{frame}

\begin{frame}{Martingale Intuition: Random Walk}

\begin{center}
\begin{tikzpicture}[scale=0.5]

% Axes
\draw[->] (-0.5,0) -- (6.5,0) node[right] {$t$};
\draw[->] (0,-3.5) -- (0,3.5) node[above] {$S_t$};

% Time ticks
\foreach \x in {1,2,3,4,5}
  \draw (\x,0.1) -- (\x,-0.1) node[below] {\small $\x$};

% Main path up to t=3
\draw[thick, BootcampBlueDark]
  (0,0) -- (1,1) -- (2,0.5) -- (3,1.2);

\node[BootcampBlueDark, left] at (3,1.2) {$S_t$};

% Possible futures from time t
\draw[thick, dashed]
  (3,1.2) -- (4,2.2);
\draw[thick, dashed]
  (3,1.2) -- (4,1.2);
\draw[thick, dashed]
  (3,1.2) -- (4,0.2);

\node at (4.1,2.2) {\small $S_t + X_{t+1}$};
\node at (4.1,1.2) {\small $S_t$};
\node at (4.1,0.2) {\small $S_t - X_{t+1}$};

% Conditional expectation
\draw[dotted, thick]
  (3,1.2) -- (4,1.2);


\end{tikzpicture}
\end{center}

\medskip

\textbf{Key message:}
\begin{itemize}
  \item Many futures are possible after time $t$.\pause
  \item Their conditional average is exactly the current value.\pause
  \item No drift $\Rightarrow$ martingale.
\end{itemize}

\end{frame}

% --------------------------------------

\begin{frame}{Example 2: Random Walk with Drift}

Let \((X_t)\) be defined by
\[
X_{t+1}=X_t+\varepsilon_{t+1},
\]
where \(\varepsilon_t\) are i.i.d.\ with mean \(m\).
\pause
\medskip

Then
\[
\mathbb{E}[X_{t+1}\mid\mathcal{F}_t]
=
X_t + m.
\]
\pause
\medskip

\begin{itemize}
  \item If \(m=0\): martingale.
  \item If \(m>0\): submartingale.
  \item If \(m<0\): supermartingale.
\end{itemize}

\end{frame}

% ======================================
\subsection{Basic Properties}

\begin{frame}{Stability Properties}

Let \((X_t)\) be a martingale.

\begin{itemize}
  \item If \(a,b\in\mathbb{R}\), then \((aX_t+b)\) is a martingale.
  \item If \((Y_t)\) is another martingale, then \((X_t+Y_t)\) is a martingale.
  \item If \(X_t\le Y_t\) a.s.\ and \((Y_t)\) is a martingale, then \((X_t)\) is a supermartingale.
\end{itemize}

\medskip

Martingales are stable under linear operations.

\end{frame}

% --------------------------------------

\begin{frame}{Jensen and Submartingales}

\begin{definitionblock}{Jensen’s Inequality for Martingales}
If \((X_t)\) is a martingale and \(g\) is convex, then\pause
\[
\mathbb{E}[g(X_{t+1})\mid\mathcal{F}_t]
\ge
g(X_t).
\]
\end{definitionblock}
\pause
\medskip

\textbf{Conclusion:}
\((g(X_t))\) is a submartingale.\pause

\medskip

\textbf{Examples:}
\begin{itemize}
  \item \(X_t^2\) is a submartingale if \(X_t\) is a martingale.\pause
  \item \(|X_t|\) is a submartingale.
\end{itemize}

\end{frame}

% ======================================
\subsection{Martingale Differences}

\begin{frame}{Martingale Differences}

Define the increments
\[
D_{t+1}=X_{t+1}-X_t.
\]
\pause
\medskip

If \((X_t)\) is a martingale, then:
\[
\mathbb{E}[D_{t+1}\mid\mathcal{F}_t]=0.
\]
\pause
\medskip

\begin{itemize}
  \item No predictable gain from the increment.\pause
  \item Martingales are sums of martingale differences.
\end{itemize}

\end{frame}

% ======================================
\subsection{Optional Sampling (Informal)}

\begin{frame}{Optional Sampling (Informal Statement)}

Let \((X_t)\) be a martingale and \(\tau\) a bounded stopping time
(with respect to the natural filtration).\pause

\medskip

Then
\[
\mathbb{E}[X_\tau]=\mathbb{E}[X_0].
\]\pause

\medskip

\begin{itemize}
  \item No strategy based on past information can create profit.\pause
  \item Formalizes the idea of a fair game.
\end{itemize}

\end{frame}

% ======================================
\subsection{Financial Interpretation}

\begin{frame}{Martingales in Finance (Discrete Time)}

Let \(S_t\) be a price process and define the discounted price
\[
\tilde S_t = \frac{S_t}{(1+r)^t}.
\]\pause

\medskip

If
\[
\mathbb{E}[\tilde S_{t+1}\mid\mathcal{F}_t]=\tilde S_t,
\]
then \((\tilde S_t)\) is a martingale.\pause

\medskip

\begin{itemize}
  \item Expected return equals risk-free rate.\pause
  \item No arbitrage using past information.
\end{itemize}

\end{frame}
% ======================================
% SECTION 4 — Applications
% ======================================
\section{Applications in Finance \& Forecasting}

\begin{frame}{Outline}
  \tableofcontents[currentsection, hideothersubsections]
\end{frame}

\subsection{Forecasting \& Risk}
\begin{frame}{Using Conditional Expectation}
\begin{itemize}
  \item Best \(L^2\)-forecast of returns given signals: \(\hat R=\mathbb{E}[R\mid \mathcal{F}_t]\).
  \item Linear factor models: projection onto span of factors.
  \item Risk attribution: \(\mathbb{E}[L\mid \mathcal{F}_t]\) and scenario conditioning.
\end{itemize}
\end{frame}

% ======================================
% SECTION 5 — Summary \& References
% ======================================

\section{References}
\begin{frame}{Suggested References}
\begin{itemize}
  \item Carl P. Simon and Lawrence Blume, \textit{Mathematics for Economists}.
  \item J. Douglas Carroll and Paul E. Green, \textit{Mathematical Tools for Applied Multivariate Analysis}.
  \item (Optional, more advanced) Bruce Hajek, \textit{Random Processes}; David Williams, \textit{Probability with Martingales}.
\end{itemize}
\end{frame}

\end{document}